\documentclass[10pt,aspectratio=169]{beamer}

% --- pdflatex-safe font & micro-typography ---
\usepackage[T1]{fontenc}
\usepackage[utf8]{inputenc}
\usepackage{lmodern}
\usepackage{microtype}

% --- math & symbols ---
\usepackage{amsmath,amssymb,mathtools,bm}

% --- hyperlinks (Beamer already loads hyperref; this is safe) ---
\hypersetup{hidelinks}

% --- theme ---
\usetheme[numbering=fraction]{metropolis}

\usepackage{pgfpages} % for notes layouts

% --- bibliography ---
\usepackage[round,authoryear]{natbib}

% --- Choose ONE of these toggles when compiling ---
% \setbeameroption{hide notes}                        % audience version (default)
% \setbeameroption{show notes}                        % notes printed under each slide
% \setbeameroption{show only notes}                   % notes-only PDF (for printing)
% \setbeameroption{show notes on second screen=right} % slide + notes side-by-side

% --- shortcuts ---
\newcommand{\E}{\mathbb{E}}
\newcommand{\R}{\mathbb{R}}
\newcommand{\btheta}{\boldsymbol{\theta}}
\newcommand{\bx}{\boldsymbol{x}}

% --- title info ---
\title{Training and generalization in overparameterized neural networks}
\subtitle{Stage 1 review}
\author{Shreyas Kalvankar}
\institute{TU Delft}
\date{\today}

\begin{document}

\maketitle

% -----------------------------------------------------------
% AGENDA SLIDE (visual blocks)
% -----------------------------------------------------------
\begin{frame}{Agenda}

	\textbf{Agenda}
	\begin{enumerate}
		\item Introduction \& Problem setup
		\item Related work: NTK, mean-field, spectral bias
		\item Research Questions \& Objectives
		\item Preliminary observations
		\item Planned work \& Discussions
	\end{enumerate}

\end{frame}
% -----------------------------------------------------------
% SECTION 1: Introduction \& Problem setup
% -----------------------------------------------------------
\begin{frame}{Introduction \& Problem setup}
	\textbf{Problem setup: Empirical Risk Minimization}
	\begin{itemize}
		\item Given a parametric model $f_\theta : \mathbb{R}^d \to \mathbb{R}$ and a dataset
		      $\{(x_i, y_i)\}_{i=1}^n$, with inputs $ x_i \in \mathbb{R}^d$ and targets $y_i \in \mathbb{R}$  (stacked as $X \in \mathbb{R}^{d \times n}, y \in \mathbb{R}^n$).
		      \pause
		\item Task: Find the parameters $\theta$
		      \pause
		\item Objective: that minimizes the empirical risk:
		      \pause
		      \only<4>{$\mathcal{L}(\theta) = \displaystyle\sum_{i=1}^{n} d(f_\theta(x_i), y_i) = \displaystyle\sum_{i=1}^{n}(f_\theta(x_i) - y_i)^2$}
		      \only<5>{$\mathcal{L}(\theta) = d(f_\theta(X), Y) = \|(f_\theta(X) - y)\|^2$}
	\end{itemize}

	Large family of these models is generally trained using algorithms like gradient descent which
	is notoriously hard to analyse.
	To motivate the obstacles to deep learning theory, we begin with a simple case.
\end{frame}

\begin{frame}{Introduction \& Problem setup}
	\textbf{Linear Networks}
	\begin{itemize}
		\item \textit{One-layer (no hidden layers):} $f_w(x) = x^\top w$, $w\in\mathbb{R}^d$
		      \pause
		\item Objective: $\mathcal{L}(w) = \frac{1}{2} \|X^\top w - y\|^2$
		\item Gradient flow: $\dot{w_t} = -\nabla_w\mathcal{L}(w_t) = -X(X^\top w_t - y) = Xy - XX^\top w_t$
		      \pause
		\item Gradient descent is a linear dynamical system governed by the data covariance matrix.
		      \pause
		\item Parameter dynamics known, can be characterized and analysed.
		      \pause
	\end{itemize}

	This is the extent to which we can do deep theoretical analysis of the training dynamics.
\end{frame}

\begin{frame}{Introduction \& Problem setup}
	\textbf{Linear Networks}
	\begin{itemize}
		\item \textit{Deep Linear (1 hidden layer of width $m$):}
		      $f_{v,W}(x) := (Wx)^\top v$ where $W \in \mathbb{R}^{m \times d}$, $v \in \mathbb{R}^{m}$.
		      \pause
		\item Objective: $\mathcal{L}(v,W) = \frac{1}{2} \|(WX)^\top v - y\|^2$
		      \pause
		\item Gradient flow:
		      \[
			      \dot v_t = -\nabla_v \mathcal{L}(v_t,W_t)
			      = W_t X \bigl(y - X^\top W_t^\top v_t\bigr),
		      \]
		      \pause
		      \[
			      \dot W_t = -\nabla_W \mathcal{L}(v_t,W_t)
			      = v_t \bigl(y - X^\top W_t^\top v_t\bigr)^\top X^\top .
		      \]
		      \pause
		\item Introducing depth already creates complex dynamics (coupled ODEs of order three) with no known general analytic
		      solutions.
	\end{itemize}

\end{frame}

\begin{frame}{Introduction \& Problem setup}
	\textbf{Non-Linear Networks}
	\begin{itemize}
		\item $f_{v,W}(X) = \phi(Wx)^\top v$ where $\phi: \mathbb{R}\to \mathbb{R}$ is a continuously differentiable (a.e.)
		      function applied elementwise, called an activation function.
		\item Objective: $\mathcal{L}(v,W) = \frac{1}{2} \|\phi(WX)^\top v - y\|^2$
		      \pause
		\item Gradient flow:
		      \[
			      \dot v_t = -\nabla_v \mathcal{L}(v_t,W_t)
			      = \phi(W_t X) \bigl(y - X^\top W_t^\top v_t\bigr),
		      \]
		      \pause
		      \[
			      \dot W_t
			      = -\nabla_W \mathcal{L}(v_t,W_t)
			      = \Big(\phi'(W_t X) \odot \big(v_t \bigl(y - \phi(W_t X)^\top v_t\bigr)^\top\big)\Big)\,X^\top.
		      \]
		      \pause
		\item Activation functions mix matrix products with elementwise
		      operations, we don't know of mathematical tools to handle these kinds of expressions.
	\end{itemize}
\end{frame}

\begin{frame}{Why Analyze Training Dynamics?}
	\begin{itemize}
		\item For linear models, gradient descent is a linear dynamical system governed by the spectrum of $XX^\top$
		      with a closed-form solution.
		      \pause
		\item For nonlinear networks, training dynamics are coupled and nonlinear, with no general closed-form description.
		      \pause
		\item But having analytical solutions allows understanding the behaviour of the system.
		      \pause
		\item This motivates studying \textbf{simplified regimes} where the dynamics become analytically tractable.
		      \pause
		\item Two prominent approaches: the \textbf{NTK} and the \textbf{mean-field} regime (we focus on the NTK).
	\end{itemize}
\end{frame}

\begin{frame}{The NTK Regime}
	\textit{Linearization at Infinite Width}

	\begin{itemize}
		\item \textbf{Motivation:} To analyze training dynamics by linearizing the network around initialization.
		\item \textbf{Formalization:}
		      \begin{itemize}
			      \item As width $m \to \infty$, the empirical kernel $\Theta_t$ converges to a deterministic, static limit $\bar\Theta$.
			            \[
				            \bar\Theta(x,x')
				            =
				            \mathbb{E}_{w}\!\big[\varphi(w^\top x)\,\varphi(w^\top x')\big]
				            +
				            \sigma_v^2\,x^\top x'\,
				            \mathbb{E}_{w}\!\big[\varphi'(w^\top x)\,\varphi'(w^\top x')\big]
			            \]
			      \item The network function $f_t$ evolves linearly with respect to this frozen kernel.
		      \end{itemize}
		\item \textbf{Training Dynamics:}
		      \begin{itemize}
			      \item Equivalent to Kernel Ridge Regression with the NTK.
			      \item Convergence is guaranteed if the limit kernel is positive definite.
		      \end{itemize}
	\end{itemize}
\end{frame}

\begin{frame}{2b. Related Work: The Mean-Field View}
	\textit{Feature Learning at Infinite Width}

	\begin{itemize}
		\item \textbf{Motivation:} To capture feature learning, which is absent in the "lazy" NTK regime.
		\item \textbf{Formalization:}
		      \begin{itemize}
			      \item Weights are treated as particles drawn from a distribution $\mu$.
			      \item Output is scaled by $1/n$ (vs $1/\sqrt{n}$ in NTK).
			            \[
				            f_{v,w}(x)
				            = \frac{1}{m} \sum_{\alpha=1}^m v_\alpha \varphi(w_\alpha^\top x)
				            =\int_\Omega v\,\varphi(w^\top x)\,d\mu(v,w)
			            \]
		      \end{itemize}
		\item \textbf{Training Dynamics:}
		      \begin{itemize}
			      \item Modeled as a Wasserstein gradient flow of the probability measure $\mu_t$.
			      \item Allows the kernel (and features) to evolve during training.
		      \end{itemize}
	\end{itemize}
\end{frame}

\begin{frame}{2c. Related Work: Spectral Bias}
	\textit{Frequency-Dependent Convergence}

	\begin{itemize}
		\item \textbf{Phenomenon:} Neural networks fit low-frequency components of the target function faster than high-frequency noise.
		\item \textbf{Theoretical Basis:}
		      \begin{itemize}
			      \item Convergence along eigen-directions is determined by the eigenvalues $\lambda_k$ of the NTK integral operator.
			      \item High frequencies correspond to small eigenvalues $\to$ slow convergence.
		      \end{itemize}
		\item \textbf{Implications:} Provides a theoretical basis for early stopping and generalization in overparameterized models.
	\end{itemize}
\end{frame}

% -----------------------------------------------------------
% SECTION 3: PRELIMINARY EXPERIMENTS
% -----------------------------------------------------------
\begin{frame}{3. Preliminary Experiments}
	\vspace{0.5cm}

	\begin{enumerate}
		\item \textbf{Function Space Convergence}
		      % Investigating width dependent effects and gradient step size stability.

		      \vspace{0.4cm}

		\item \textbf{Kernel Drift \& Regime Transitions}
		      % Distinguishing between feature learning (Mean-Field) and linear (NTK) regimes.

		      \vspace{0.4cm}

		\item \textbf{Spectral Analysis of Empirical NTK}
		      % Eigendecomposition to verify spectral bias on real data.
	\end{enumerate}
\end{frame}

% -----------------------------------------------------------
% CONCLUSION
% -----------------------------------------------------------
\begin{frame}{Summary \& Discussion}
	\textbf{Structure of Introduction \& Literature Review:}
	\begin{enumerate}
		\item \textbf{Foundations:} From linear regression to deep linear dynamics.
		\item \textbf{Infinite Limits:} NTK limit (static kernel), Mean-Field limit (changing kernel).
		\item \textbf{Spectral Properties:} How eigenvalues dictate learnability (Spectral Bias).
	\end{enumerate}

	\vspace{0.5cm}
	\textit{Question for Supervisors: Does this structure logically support the move to finite-width deviations in later chapters?}
\end{frame}

% -----------------------
\end{document}
